% ============================================================================
% CHAPTER 1 SOLUTIONS GUIDE
% Exercise Solutions & Answer Keys
% ============================================================================
% Source: /markdown/ch1/C1AM_updated.md (Solutions Guide)
%         /markdown/ch1/Ch1AS_updated.md (Technical Exercise Solutions)
% Note: This file is for the Instructor's Edition, not the main book
% ============================================================================

\chapter{Solutions Guide: Chapter 1}
\label{ch:solutions-ch01}

\begin{chapterquote}{Complete solutions for all exercises, activities, and discussion questions}
Framework-aligned answers enabling instructors to facilitate discussions with confidence and consistency.
\end{chapterquote}

% ============================================================================
\section{Discussion Question Solutions}
% ============================================================================

\subsection{Introductory Level Solutions}

\subsubsection{Question 1: Recognition Examples}

\textbf{Prompt:} ``Give three examples of technology asking you for help. Why do you think it asks instead of just making a decision?''

\paragraph{Sample Strong Answers.}

\begin{enumerate}
    \item \textbf{Banking app asking to confirm unusual purchases}
        \begin{itemize}
            \item \emph{Why:} High cost of false positives (blocking legitimate transactions) vs.\ low cost of asking
            \item \emph{Uncertainty type:} Unusual pattern detection with high stakes
            \item \emph{Framework dimension:} Stakes Calibration---financial impact justifies interruption
        \end{itemize}

    \item \textbf{Voice assistant saying ``I found this on the web''}
        \begin{itemize}
            \item \emph{Why:} Low confidence in ability to directly answer; safer to provide source than guess
            \item \emph{Uncertainty type:} Knowledge boundary recognition
            \item \emph{Framework dimension:} Uncertainty Detection---recognizes it doesn't ``know'' the answer
        \end{itemize}

    \item \textbf{Photo app asking ``Is this [person's name]?''}
        \begin{itemize}
            \item \emph{Why:} Facial recognition confidence below threshold; privacy implications of wrong tagging
            \item \emph{Uncertainty type:} Aleatoric uncertainty due to image quality/angle
            \item \emph{Framework dimension:} Intervention Design---simple yes/no reduces cognitive load
        \end{itemize}
\end{enumerate}

\paragraph{Grading Criteria.}
\begin{description}
    \item[Excellent] Three clear examples with accurate analysis of why the system asks, references to framework dimensions
    \item[Good] Three examples with basic understanding of uncertainty concepts
    \item[Satisfactory] Three examples but limited analysis
    \item[Needs Improvement] Fewer than three examples or misunderstands the concept
\end{description}

\subsubsection{Question 2: Netflix Analysis}

\textbf{Prompt:} ``Netflix could easily keep playing episodes forever. Why doesn't it? What are the costs and benefits?''

\paragraph{Model Answer.}

\textbf{Costs of asking:}
\begin{itemize}
    \item User annoyance/interruption
    \item Potential for users to stop watching when they would have continued
    \item Development and maintenance of detection algorithms
\end{itemize}

\textbf{Benefits of asking:}
\begin{itemize}
    \item Bandwidth savings (significant operational cost)
    \item Better user experience (don't lose place in series)
    \item Reduced server load
    \item More accurate viewing analytics
    \item Regulatory compliance (data usage transparency)
\end{itemize}

\textbf{Framework Analysis:}
\begin{itemize}
    \item \textbf{Uncertainty Detection:} Netflix monitors engagement signals (interaction, time)
    \item \textbf{Timing:} After 3 episodes or 90 minutes---rules-based, predictable
    \item \textbf{Stakes Calibration:} Low stakes (minor annoyance) justify asking even when uncertain
    \item \textbf{Intervention Design:} Simple yes/no---minimal cognitive load
    \item \textbf{Feedback Integration:} User response informs future thresholds
\end{itemize}

\begin{keyconcept}[Why They Ask]
Cost-benefit analysis shows savings outweigh interruption costs. The Stakes Calibration dimension tells us: when stakes are low and uncertainty is moderate, asking is almost always the right choice.
\end{keyconcept}

\paragraph{Student Misconceptions to Address.}
\begin{itemize}
    \item ``It's just about bandwidth'' (ignores user experience benefits)
    \item ``They should detect activity differently'' (explain current detection limitations)
    \item ``It's always annoying'' (help students recognize when they appreciate it)
\end{itemize}

\subsubsection{Question 3: Netflix vs.\ Nest Comparison}

\textbf{Prompt:} ``Compare the Netflix `Are you still watching?' prompt to the Nest thermostat's Home/Away detection. One asks for help, one doesn't. What are the consequences of each approach?''

\paragraph{Model Answer.}

\begin{table}[htbp]
\caption{Netflix vs.\ Nest: five-dimension comparison}
\label{tab:netflix-nest-five-dim}
\centering
\begin{tabular}{@{}p{3.5cm}p{4.5cm}p{4.5cm}@{}}
\toprule
\textbf{Aspect} & \textbf{Netflix} & \textbf{Nest Thermostat} \\
\midrule
Uncertainty Detection & Monitors engagement signals & Binary decision from sensors \\
Intervention Design & Simple yes/no button & No intervention mechanism \\
Timing & Clear rules (3 episodes/90 min) & Variable (30 min to 2 hours) \\
Stakes Calibration & Low stakes, can afford to ask & Doesn't distinguish mild discomfort from pipes freezing \\
Feedback Integration & Learns from responses & Limited---requires manual override \\
\bottomrule
\end{tabular}
\end{table}

\paragraph{Consequences.}

\textbf{Netflix (asks):}
\begin{itemize}
    \item Occasional minor annoyance
    \item Saves bandwidth when users are actually away
    \item Users feel respected (``it asked instead of assuming'')
    \item Rare failures (interrupts engaged user) are low-cost
\end{itemize}

\textbf{Nest (doesn't ask):}
\begin{itemize}
    \item Users wake up freezing: ``I woke up to a cold house because it thought I was away''
    \item Forum complaints: ``Nest keeps saying I'm away when I'm home''
    \item Trust erosion: Users disable ``smart'' features, defeating the purpose
    \item High-cost failures when the system is wrong
\end{itemize}

\begin{keyconcept}[Key Insight]
The system that asks is perceived as ``smarter'' even though it might seem less autonomous. The system that acts confidently without asking creates more user frustration.
\end{keyconcept}

\subsubsection{Question 4: Nest Thermostat Redesign}

\textbf{Prompt:} ``If you were redesigning the Nest thermostat, when would you want it to ask the human for input versus making decisions automatically?''

\paragraph{Sample Solution Framework.}

\textbf{Automatic Decisions (High Confidence):}
\begin{itemize}
    \item Small temperature adjustments ($\pm 2$°F) within comfortable range
    \item Schedule-based changes when pattern is well-established
    \item Energy-saving mode during confirmed long absences (e.g., vacation mode)
    \item Weather-based preemptive adjustments (heating before cold front)
\end{itemize}

\textbf{Ask for Input (Uncertainty Scenarios):}
\begin{itemize}
    \item First detection of ``Away'' after long presence period
    \item Unusual occupancy patterns detected
    \item Extreme outside temperature conditions (higher stakes)
    \item Conflicting signals (motion but no phone, or phone but no motion)
    \item First-time seasonal transitions
    \item Manual override patterns suggest dissatisfaction
\end{itemize}

\paragraph{Implementation Strategy (Five-Dimension Framework).}

\begin{lstlisting}[style=python, caption={Nest redesign decision logic}]
def should_ask_or_act(presence_confidence, stakes, recent_overrides):
    # Stakes Calibration
    if extreme_weather():
        required_confidence = 0.95  # High stakes = need high confidence
    else:
        required_confidence = 0.70  # Lower stakes = can act with less certainty

    # Uncertainty Detection + Feedback Integration
    if recent_overrides > 2:
        required_confidence += 0.15  # User has been correcting us

    # Intervention Decision
    if presence_confidence > required_confidence:
        return "act_automatically"
    elif presence_confidence < (1 - required_confidence):
        return "act_automatically"  # Very confident they're away
    else:
        return "ask_user"
\end{lstlisting}

\paragraph{Interface Design (Intervention Design dimension).}
\begin{verbatim}
Notification: "Haven't seen activity for 2 hours.
Should I switch to energy-saving mode?
[Yes, I'm out] [No, I'm home] [Ask me in 1 hour]"
\end{verbatim}

\subsection{Intermediate Level Solutions}

\subsubsection{Framework Application: Uber Surge Pricing}

\textbf{Prompt:} ``Apply the five-dimension framework to analyze Uber's surge pricing notifications.''

\paragraph{Sample Strong Answer.}

\begin{table}[htbp]
\caption{Uber surge pricing: five-dimension analysis}
\label{tab:uber-surge}
\centering
\begin{tabular}{@{}p{3.5cm}p{1.5cm}p{7cm}@{}}
\toprule
\textbf{Dimension} & \textbf{Rating} & \textbf{Analysis} \\
\midrule
Uncertainty Detection & 4/5 & System knows demand/supply imbalance with high confidence; uncertainty is in user's willingness to pay \\
Intervention Design & 4/5 & Clear notification with exact multiplier, requires explicit acceptance (``Accept 2.3x surge'') \\
Timing & 5/5 & Asks at exactly the right moment---before booking, when decision matters \\
Stakes Calibration & 3/5 & Treats all rides similarly; doesn't distinguish urgent medical need from casual trip \\
Feedback Integration & 4/5 & User acceptance/rejection directly affects future pricing algorithms \\
\bottomrule
\end{tabular}
\end{table}

\textbf{Weakest Dimension:} Stakes Calibration. The system doesn't know \emph{why} you need a ride. A 3$\times$ surge for a bar-to-home trip is different from 3$\times$ surge for a hospital emergency.

\textbf{Improvement Suggestion:} Allow users to indicate ``urgent'' rides that bypass surge (perhaps with post-ride verification) or at least provide context (``High demand due to rain'' helps users decide).

\subsubsection{Design Question: AI Essay Grading}

\textbf{Prompt:} ``Design the `asking for help' strategy for an AI system that helps teachers grade student essays.''

\paragraph{Comprehensive Solution.}

\textbf{Uncertainty Triggers (Uncertainty Detection):}
\begin{enumerate}
    \item \textbf{Content Analysis Uncertainty:}
        \begin{itemize}
            \item Unusual writing style for known student (possible plagiarism or growth)
            \item Topic drift or unclear thesis
            \item Creative/artistic responses to structured prompts
            \item Potential plagiarism flags with low confidence
        \end{itemize}

    \item \textbf{Technical Quality Uncertainty:}
        \begin{itemize}
            \item Borderline grammar/spelling issues
            \item Complex sentence structures difficult to parse
            \item Non-standard English dialects or ESL patterns
        \end{itemize}

    \item \textbf{Subjective Assessment Uncertainty:}
        \begin{itemize}
            \item Responses near grade boundaries (B+/A- borderline)
            \item Creative interpretations of assignment requirements
            \item Evidence of exceptional insight that doesn't fit rubric
        \end{itemize}
\end{enumerate}

\textbf{Human Intervention Strategy (Timing + Stakes Calibration):}
\begin{verbatim}
High Priority (Immediate Review):
  - Potential academic integrity violations (high stakes)
  - Grades >1 letter different from student's historical average
  - Responses AI cannot categorize at all

Medium Priority (Batch Review):
  - Borderline grades (within 5 points of cutoff)
  - Technical issues affecting readability
  - Unusual but potentially valuable creative responses

Low Priority (Optional Review):
  - Random sample for calibration (maintains system accuracy)
  - Review AI reasoning for transparent grading
\end{verbatim}

\textbf{Interface Design (Intervention Design):}
\begin{itemize}
    \item Show AI confidence scores to teacher
    \item Highlight specific sentences/phrases causing uncertainty
    \item Provide suggested grade \emph{ranges} rather than point values
    \item Include links to similar essays for comparison
    \item Offer ``quick confirm'' vs.\ ``full review'' options
\end{itemize}

\textbf{Feedback Integration:}
\begin{itemize}
    \item Teacher corrections update model understanding
    \item Track which types of essays cause most uncertainty
    \item Adjust confidence thresholds based on teacher override patterns
\end{itemize}

\subsection{Advanced Level Solutions}

\subsubsection{Technical Implementation Question}

\textbf{Prompt:} ``How would you implement an optimal intervention threshold for a medical diagnosis system?''

\paragraph{Solution Framework.}

\textbf{1. Data Collection:}
\begin{lstlisting}[style=python, caption={Medical AI data collection setup}]
training_data = {
    'medical_images': preprocessed_scans,
    'expert_diagnoses': ground_truth_labels,
    'confidence_scores': model_uncertainty_estimates,
    'human_review_outcomes': expert_corrections,
    'cost_parameters': {
        'false_positive_cost': delayed_treatment_cost,  # ~$200
        'false_negative_cost': missed_diagnosis_cost,   # ~$50,000+
        'human_review_cost': radiologist_time_cost      # ~$50
    }
}
\end{lstlisting}

\textbf{2. Stakes-Calibrated Threshold:}
\begin{lstlisting}[style=python, caption={Optimal threshold calculation}]
def calculate_optimal_threshold(costs):
    """
    Using the Stakes Calibration dimension:
    theta* = C_FP / (C_FP + C_FN)

    For medical diagnosis with high FN costs:
    theta* = 200 / (200 + 50000) = 0.004

    This is the positive-class decision threshold: classify as disease-present
    (and send for treatment/follow-up) whenever P(disease) > 0.004.
    Equivalently, only auto-dismiss as healthy when P(no disease) > 0.996,
    i.e., when confidence in the negative class exceeds 99.6%.
    In a HITL review zone: flag any case where P(no disease) < 99.6%.
    """
    threshold = costs['false_positive_cost'] / (
        costs['false_positive_cost'] + costs['false_negative_cost']
    )
    return threshold
\end{lstlisting}

\textbf{3. Uncertainty Estimation:}
\begin{lstlisting}[style=python, caption={Ensemble and MC Dropout uncertainty estimation}]
def estimate_uncertainty(model, input_image, method='ensemble'):
    if method == 'ensemble':
        predictions = [model_i(input_image) for model_i in ensemble]
        mean_pred = np.mean(predictions, axis=0)
        uncertainty = np.std(predictions, axis=0)

    elif method == 'mc_dropout':
        predictions = []
        for _ in range(n_samples):
            model.train()  # Enable dropout
            pred = model(input_image)
            predictions.append(pred)

        mean_pred = torch.mean(torch.stack(predictions), dim=0)
        uncertainty = torch.std(torch.stack(predictions), dim=0)

    return mean_pred, uncertainty
\end{lstlisting}

\textbf{4. Evaluation Metrics (Five-Dimension Assessment):}
\begin{lstlisting}[style=python, caption={HITL system evaluation}]
def evaluate_hitl_system(threshold, test_data):
    metrics = {
        # Uncertainty Detection quality
        'calibration_error': expected_calibration_error(
            confidences, accuracies),

        # Intervention quality
        'precision_intervention': (tp_interventions /
            (tp_interventions + fp_interventions)),
        'recall_intervention': (tp_interventions /
            (tp_interventions + fn_interventions)),

        # System-level performance
        'system_accuracy': correct_predictions / total_predictions,

        # Cost-effectiveness (Stakes Calibration)
        'cost_effectiveness': total_benefit / total_cost,

        # Human factors
        'human_workload': interventions_per_hour,
    }
    return metrics
\end{lstlisting}

% ============================================================================
\section{Activity Solutions}
% ============================================================================

\subsection{Activity 1: HITL Detection Hunt---Solution Guide}

\textbf{Setup:} Students find 5 examples of systems asking for human input.

\paragraph{Sample Student Findings with Framework Analysis.}

\subsubsection{Strong Example: Two-Factor Authentication}
\begin{description}
    \item[System] Banking/email login
    \item[Uncertainty Detection] Unusual login location/device detected
    \item[Intervention Design] Code sent via SMS, simple entry
    \item[Timing] At moment of login attempt (synchronous)
    \item[Stakes Calibration] High (security breach) justifies interruption
    \item[Feedback Integration] Successful logins from new device update trust
\end{description}

\subsubsection{Strong Example: ``Did you mean?'' Search Suggestions}
\begin{description}
    \item[System] Google/search engines
    \item[Uncertainty Detection] Typo probability calculated
    \item[Intervention Design] Shows both original and suggestions (low commitment)
    \item[Timing] After search, before results (perfect timing)
    \item[Stakes Calibration] Low stakes, can afford to suggest
    \item[Feedback Integration] Click on suggestion = confirmation
\end{description}

\paragraph{Framework Classification Exercise.}

Have students sort their findings by weakest dimension:

\begin{table}[htbp]
\caption{Example systems sorted by weakest dimension}
\label{tab:weakest-dimension}
\centering
\begin{tabular}{@{}p{4cm}p{3.5cm}p{4cm}@{}}
\toprule
\textbf{Example} & \textbf{Weakest Dimension} & \textbf{Why} \\
\midrule
Cookie consent banners & Timing & Asks before user engagement \\
App permission requests & Stakes Calibration & Treats all permissions equally \\
Voice assistant ``I don't understand'' & Intervention Design & Doesn't explain \emph{why} confused \\
\bottomrule
\end{tabular}
\end{table}

\subsection{Activity 2: Nest Thermostat Redesign---Solution Guide}

\textbf{Scenario:} Redesign Nest Home/Away detection using the five-dimension framework.

\paragraph{Problem Analysis.}

Current Nest fails on multiple dimensions:
\begin{itemize}
    \item \textbf{Uncertainty Detection:} Binary home/away, no confidence expression
    \item \textbf{Intervention Design:} Never asks
    \item \textbf{Stakes Calibration:} Doesn't adjust for weather conditions
    \item \textbf{Feedback Integration:} Manual overrides don't improve future predictions
\end{itemize}

\paragraph{Redesign Solution.}

\textbf{1. Uncertainty Detection Enhancement:}
\begin{lstlisting}[style=python, caption={Improved presence confidence calculation}]
def calculate_presence_confidence(signals):
    confidence_score = 0.5  # Start neutral

    if motion_detected_recently:
        confidence_score += 0.3
    if phone_in_geofence:
        confidence_score += 0.25
    if typical_home_time:
        confidence_score += 0.15
    if smart_device_activity:  # TV, lights, etc.
        confidence_score += 0.1

    return min(confidence_score, 1.0)
\end{lstlisting}

\textbf{2. Stakes-Adjusted Thresholds:}
\begin{lstlisting}[style=python, caption={Stakes-adjusted presence threshold}]
def get_threshold(weather, current_setting):
    base_threshold = 0.7

    # Higher stakes = need more confidence
    if outside_temp < 32:  # Freezing
        return 0.95
    elif outside_temp > 90:  # Very hot
        return 0.90
    elif abs(outside_temp - current_setting) > 20:
        return 0.85
    else:
        return base_threshold
\end{lstlisting}

\textbf{3. Intervention Design:}
\begin{verbatim}
+------------------------------------------+
|  Nest Home Status                        |
|                                          |
|  I haven't detected activity for         |
|  2 hours. It's 28 degrees F outside.     |
|                                          |
|  Should I switch to Away mode?           |
|                                          |
|  [Yes, I'm out]  [No, I'm home]         |
|  [Snooze 1 hour] [Don't ask today]      |
+------------------------------------------+
\end{verbatim}

\textbf{4. Feedback Integration:}
\begin{lstlisting}[style=python, caption={Feedback integration for Nest redesign}]
def update_model(user_response, context):
    if response == "no_im_home":
        # We were wrong - lower confidence in similar situations
        adjust_presence_model(context, direction="more_conservative")

    if response == "dont_ask_today":
        # User is working from home - learn this pattern
        learn_wfh_pattern(context.day_of_week)
\end{lstlisting}

\subsection{Activity 3: Cost-Benefit Analysis---Solution Guide}

\textbf{Scenario:} Email spam filter intervention frequency.

\paragraph{Given Parameters.}
\begin{itemize}
    \item Spam detection accuracy: 95\% without human intervention
    \item Human review accuracy: 99\%
    \item Cost of missed spam: \$2 (user time + annoyance)
    \item Cost of blocked legitimate email: \$50 (business impact)
    \item Cost of human review: \$0.10 per email
    \item Spam rate: 40\%
\end{itemize}

\paragraph{Mathematical Solution.}

\textbf{Without Human Review:}
\[
\text{Expected cost per email}
= P(\text{spam}) \cdot P(\text{miss} \mid \text{spam}) \cdot C_{\text{miss}}
+ P(\lnot\text{spam}) \cdot P(\text{block} \mid \lnot\text{spam}) \cdot C_{\text{block}}
\]
\[
= 0.40 \times 0.05 \times \$2 + 0.60 \times 0.05 \times \$50
= \$0.04 + \$1.50 = \$1.54 \text{ per email}
\]

\textbf{With Optimal Human Review:}
\begin{lstlisting}[style=python, caption={Optimal spam filter threshold calculation}]
def calculate_optimal_threshold():
    best_cost = float('inf')
    best_threshold = 0

    for review_rate in np.linspace(0, 1, 100):
        # Assume we review emails where confidence is lowest
        # Review improves accuracy to 99%

        auto_cost = (1 - review_rate) * 1.54  # Unreviewed emails
        review_cost = review_rate * (
            0.01 * 0.40 * 2 +      # Reviewed, still miss spam
            0.01 * 0.60 * 50 +     # Reviewed, still block legitimate
            0.10                    # Review cost
        )

        total_cost = auto_cost + review_cost

        if total_cost < best_cost:
            best_cost = total_cost
            best_threshold = review_rate

    return best_threshold, best_cost

# Result: Review ~30% of emails (lowest confidence ones)
# Expected cost drops from $1.54 to ~$0.95 per email
\end{lstlisting}

\begin{keyconcept}[Key Insight]
Stakes Calibration matters---the high cost of blocking legitimate email (\$50 vs.\ \$2 for spam) means we should be more aggressive about reviewing emails the system thinks are spam.
\end{keyconcept}

% ============================================================================
\section{``Try This'' Exercise---Sample Student Responses}
% ============================================================================

\textbf{Assignment:} Count HITL moments over 24 hours.

\paragraph{Sample Strong Response.}

\begin{storyhook}
``I counted 23 HITL moments in 24 hours. Here are the most interesting:

\textbf{1. Bank app fraud alert} (morning coffee purchase at new location)
\begin{itemize}
    \item Timing: Perfect---immediate, while I still had context
    \item Intervention: Simple ``Yes/No'' text response
    \item Would improve: Tell me \emph{why} it flagged (was it the location or amount?)
\end{itemize}

\textbf{2. Google Maps ``Faster route available''} (during commute)
\begin{itemize}
    \item Timing: Good---offered alternative without forcing change
    \item Stakes awareness: High traffic = more aggressive about suggesting
    \item Weakness: Didn't explain \emph{how much} faster (2 min vs.\ 20 min matters)
\end{itemize}

\textbf{3. iPhone asking to update} (6th time this week)
\begin{itemize}
    \item Timing: TERRIBLE---asked during video call
    \item Clear alert fatigue---I've started ignoring these
    \item Framework failure: No Stakes Calibration (treats all moments equally)
\end{itemize}

\textbf{4. Nest thermostat switched to Away} (while I was reading quietly)
\begin{itemize}
    \item This is the example from the chapter! It DIDN'T ask.
    \item I was home but hadn't walked past the sensor.
    \item Came back to a 62°F house.
\end{itemize}

\textbf{Patterns I noticed:}
\begin{itemize}
    \item Systems that ask are generally more trustworthy
    \item Financial apps have the best intervention timing
    \item Smart home devices have the worst (or no) intervention design
    \item The more a system asks unnecessarily, the more I ignore it
\end{itemize}''
\end{storyhook}

% ============================================================================
\section{Quick Reference: Common Student Questions \& Answers}
% ============================================================================

\subsection{``Why don't they just make the AI better instead of asking humans?''}

\textbf{Answer:} There are fundamental limits to what any system can know with certainty. Even perfect AI will encounter:
\begin{itemize}
    \item Situations not in training data (epistemic uncertainty)
    \item Inherently random events (aleatoric uncertainty)
    \item Changing environments and user preferences
    \item Edge cases where stakes are too high for any uncertainty
\end{itemize}

\textbf{Framework connection:} The Uncertainty Detection dimension isn't about eliminating uncertainty---it's about \emph{recognizing} it. Even a perfect AI would still need to know when to ask.

\subsection{``Isn't asking just an excuse for lazy programming?''}

\textbf{Answer:} Actually the opposite---implementing good uncertainty quantification requires sophisticated engineering:
\begin{itemize}
    \item Confidence estimation algorithms
    \item Calibration methods
    \item User interface design (Intervention Design)
    \item Cost-benefit optimization (Stakes Calibration)
    \item A/B testing frameworks
    \item Feedback integration systems
\end{itemize}

\textbf{Teaching moment:} Show the Technical Appendix complexity. Asking well is harder than not asking at all.

\subsection{``How do you know when the human is wrong?''}

\textbf{Answer:}
\begin{itemize}
    \item Humans aren't always right, but they're right about \emph{different things} than AI
    \item Multiple human opinions can be aggregated
    \item Domain expertise matters (radiologist vs.\ random person)
    \item Systems can learn from patterns in human corrections (Feedback Integration)
    \item The goal is better \emph{joint} performance, not perfect individual performance
\end{itemize}

\textbf{Framework connection:} The Feedback Integration dimension handles this---good systems learn from human responses over time, even when humans are occasionally wrong.

\subsection{``Why did the Air Canada chatbot case matter legally?''}

\textbf{Answer:}
\begin{itemize}
    \item Tribunal explicitly rejected the ``chatbot is a separate entity'' argument
    \item Companies are responsible for ALL information on their website, including AI-generated content
    \item Sets precedent: You can't blame your tools for your failures
    \item Implies companies MUST build uncertainty detection into customer-facing AI
\end{itemize}

\begin{cautionbox}[Framework Connection]
Air Canada's chatbot failed on Uncertainty Detection (didn't know it was wrong) AND Intervention Design (no escalation path). This wasn't just bad UX---it was legally liable.
\end{cautionbox}

% ============================================================================
\section{Assessment Solution Keys}
% ============================================================================

\subsection{Framework Analysis Paper---Sample A-Grade Response}

\textbf{Technology Analyzed:} Google Maps Navigation

\begin{table}[htbp]
\caption{Google Maps: five-dimension framework analysis (sample A-grade)}
\label{tab:google-maps-analysis}
\centering
\begin{tabular}{@{}p{3.5cm}p{1.5cm}p{7cm}@{}}
\toprule
\textbf{Dimension} & \textbf{Rating} & \textbf{Evidence} \\
\midrule
Uncertainty Detection & 5/5 & Displays confidence in ETA as range (``15--22 min''), acknowledges traffic uncertainty \\
Intervention Design & 4/5 & Offers alternatives clearly, but sometimes too many options overwhelm \\
Timing & 5/5 & Re-routing suggestions come at decision points (intersections), not mid-highway \\
Stakes Calibration & 3/5 & Doesn't distinguish ``running late for flight'' from ``casual drive'' \\
Feedback Integration & 5/5 & User route choices improve future suggestions; traffic reports integrated \\
\bottomrule
\end{tabular}
\end{table}

\paragraph{Weakest Dimension Analysis.}

Stakes Calibration is the weakest dimension. Google Maps treats all trips equally---a 5-minute delay is presented the same way whether I'm leisurely exploring or racing to catch a flight. This could be improved by:
\begin{enumerate}
    \item Allowing users to mark ``time-sensitive'' trips
    \item Learning from calendar integration (meeting in 20 min = higher stakes)
    \item Adjusting notification urgency based on buffer time
\end{enumerate}

\paragraph{Improvement Proposal.}

Add a ``Priority Mode'' that:
\begin{itemize}
    \item Asks users to indicate urgency before starting navigation
    \item Adjusts re-routing aggressiveness based on stakes
    \item Provides proactive delay warnings for time-sensitive trips
    \item Learns which trips are typically time-sensitive (work commute vs.\ weekend drive)
\end{itemize}

This would strengthen the Stakes Calibration dimension without sacrificing the strong Uncertainty Detection and Timing that already exist.

% ============================================================================
% TECHNICAL EXERCISE SOLUTIONS
% ============================================================================

\section{Technical Exercise Solutions}
\label{sec:technical-exercise-solutions}

\subsection{Exercise 1: Five-Dimension Framework Analysis}

\textbf{Task:} Apply the five-dimension framework to analyze three systems from your daily life. For each system, rate each dimension (1--5) and identify the weakest dimension.

\subsubsection{System 1: Spotify ``Discover Weekly'' Recommendations}

\begin{table}[htbp]
\caption{Spotify Discover Weekly: five-dimension analysis}
\label{tab:spotify-analysis}
\centering
\begin{tabular}{@{}p{3.5cm}p{1.5cm}p{7cm}@{}}
\toprule
\textbf{Dimension} & \textbf{Rating} & \textbf{Evidence} \\
\midrule
Uncertainty Detection & 3/5 & Doesn't explicitly show confidence; mixes high and low confidence recommendations without distinction \\
Intervention Design & 4/5 & Like/dislike buttons, ``Don't play this'' option, easy to skip \\
Timing & 4/5 & Weekly refresh is predictable; doesn't interrupt current listening \\
Stakes Calibration & 5/5 & Low stakes (just music), can afford to experiment \\
Feedback Integration & 5/5 & Strong learning from likes, skips, and listening time \\
\bottomrule
\end{tabular}
\end{table}

\textbf{Weakest Dimension:} Uncertainty Detection.

\textbf{Analysis:} Spotify doesn't distinguish between ``I'm confident you'll love this'' and ``This is experimental.'' Users can't tell which recommendations are safe bets vs.\ explorations. Improvement: Add confidence indicators or separate ``Confident Picks'' from ``Experiments.''

\subsubsection{System 2: iPhone Face ID Unlock}

\begin{table}[htbp]
\caption{iPhone Face ID: five-dimension analysis}
\label{tab:faceid-analysis}
\centering
\begin{tabular}{@{}p{3.5cm}p{1.5cm}p{7cm}@{}}
\toprule
\textbf{Dimension} & \textbf{Rating} & \textbf{Evidence} \\
\midrule
Uncertainty Detection & 5/5 & Clear confidence threshold---either unlocks or doesn't \\
Intervention Design & 4/5 & Falls back to passcode when uncertain; clear flow \\
Timing & 5/5 & Instant attempt on wake, immediate fallback \\
Stakes Calibration & 4/5 & Security-conscious (won't unlock if unsure), but treats all unlocks equally \\
Feedback Integration & 4/5 & Learns face changes over time, but doesn't explain learning \\
\bottomrule
\end{tabular}
\end{table}

\textbf{Weakest Dimension:} Stakes Calibration (marginal).

\textbf{Analysis:} Face ID treats unlocking to check the time the same as unlocking to access a banking app. Could adjust confidence thresholds based on what is being accessed.

\subsubsection{System 3: Gmail Smart Compose}

\begin{table}[htbp]
\caption{Gmail Smart Compose: five-dimension analysis}
\label{tab:gmail-analysis}
\centering
\begin{tabular}{@{}p{3.5cm}p{1.5cm}p{7cm}@{}}
\toprule
\textbf{Dimension} & \textbf{Rating} & \textbf{Evidence} \\
\midrule
Uncertainty Detection & 4/5 & Only suggests when confident; grey text indicates suggestion \\
Intervention Design & 5/5 & Tab to accept, ignore to reject---minimal friction \\
Timing & 5/5 & Appears at natural pause points in typing \\
Stakes Calibration & 2/5 & Doesn't know if email is to boss vs.\ friend \\
Feedback Integration & 3/5 & Learns from accepts/rejects but slowly \\
\bottomrule
\end{tabular}
\end{table}

\textbf{Weakest Dimension:} Stakes Calibration.

\textbf{Analysis:} Smart Compose suggests casual completions (``Thanks!'' ``Sounds good!'') in professional contexts. Doesn't distinguish formal vs.\ informal emails. Could detect recipient (manager vs.\ friend) and adjust formality.

\subsection{Exercise 2: Uncertainty Estimation Implementation}

\textbf{Task:} Implement both Monte Carlo Dropout and Deep Ensembles for a simple classification task. Compare their uncertainty estimates and computational costs.

\begin{lstlisting}[style=python, caption={Complete uncertainty estimation implementation}]
import torch
import torch.nn as nn
import torch.nn.functional as F
import numpy as np
import matplotlib.pyplot as plt
from sklearn.datasets import make_classification
from sklearn.model_selection import train_test_split
import time

# ============================================
# SETUP: Generate Dataset
# ============================================

X, y = make_classification(
    n_samples=2000,
    n_features=20,
    n_classes=2,
    n_informative=10,
    n_redundant=5,
    random_state=42
)

X_train, X_test, y_train, y_test = train_test_split(
    X, y, test_size=0.3, random_state=42
)

X_train = torch.FloatTensor(X_train)
X_test = torch.FloatTensor(X_test)
y_train = torch.LongTensor(y_train)
y_test = torch.LongTensor(y_test)

# ============================================
# MODEL: Classifier with Dropout
# ============================================

class DropoutClassifier(nn.Module):
    """
    Simple classifier with dropout layers for MC Dropout
    uncertainty estimation.
    """
    def __init__(self, input_dim, dropout_rate=0.2):
        super(DropoutClassifier, self).__init__()
        self.fc1 = nn.Linear(input_dim, 64)
        self.fc2 = nn.Linear(64, 32)
        self.fc3 = nn.Linear(32, 2)
        self.dropout = nn.Dropout(dropout_rate)

    def forward(self, x):
        x = F.relu(self.fc1(x))
        x = self.dropout(x)
        x = F.relu(self.fc2(x))
        x = self.dropout(x)
        x = self.fc3(x)
        return F.softmax(x, dim=1)

# ============================================
# METHOD 1: Monte Carlo Dropout
# ============================================

class MCDropoutPredictor:
    """
    Monte Carlo Dropout for uncertainty estimation.

    Key insight: By keeping dropout enabled during inference and running
    multiple forward passes, we get a distribution of predictions that
    reflects model uncertainty.
    """
    def __init__(self, model, n_samples=100):
        self.model = model
        self.n_samples = n_samples

    def predict_with_uncertainty(self, x):
        """Run multiple forward passes with dropout enabled."""
        self.model.train()  # Keep dropout active!
        predictions = []

        start_time = time.time()
        with torch.no_grad():
            for _ in range(self.n_samples):
                pred = self.model(x)
                predictions.append(pred.numpy())

        predictions = np.array(predictions)
        mean_pred = np.mean(predictions, axis=0)
        uncertainty = np.std(predictions, axis=0)

        inference_time = time.time() - start_time
        return mean_pred, uncertainty, inference_time

# ============================================
# METHOD 2: Deep Ensembles
# ============================================

class DeepEnsemble:
    """
    Deep Ensemble for uncertainty estimation.

    Key insight: Train multiple models with different initializations.
    Disagreement between models indicates epistemic uncertainty.
    """
    def __init__(self, input_dim, n_models=5):
        self.models = []
        self.n_models = n_models

        for i in range(n_models):
            torch.manual_seed(i * 42)
            model = DropoutClassifier(input_dim, dropout_rate=0.0)
            self.models.append(model)

    def train_ensemble(self, X_train, y_train, epochs=100, verbose=True):
        for i, model in enumerate(self.models):
            if verbose:
                print(f"Training model {i+1}/{self.n_models}...")

            optimizer = torch.optim.Adam(model.parameters(), lr=0.001)
            criterion = nn.CrossEntropyLoss()

            for epoch in range(epochs):
                optimizer.zero_grad()
                outputs = model(X_train)
                loss = criterion(outputs, y_train)
                loss.backward()
                optimizer.step()

    def predict_with_uncertainty(self, x):
        predictions = []

        start_time = time.time()
        for model in self.models:
            model.eval()
            with torch.no_grad():
                pred = model(x)
                predictions.append(pred.numpy())

        predictions = np.array(predictions)
        mean_pred = np.mean(predictions, axis=0)
        epistemic_uncertainty = np.std(predictions, axis=0)

        inference_time = time.time() - start_time
        return mean_pred, epistemic_uncertainty, inference_time

# ============================================
# TRAINING AND COMPARISON
# ============================================

# Train MC Dropout model
mc_model = DropoutClassifier(X_train.shape[1])
optimizer = torch.optim.Adam(mc_model.parameters(), lr=0.001)
criterion = nn.CrossEntropyLoss()

for epoch in range(100):
    optimizer.zero_grad()
    outputs = mc_model(X_train)
    loss = criterion(outputs, y_train)
    loss.backward()
    optimizer.step()

# Train Deep Ensemble
ensemble = DeepEnsemble(X_train.shape[1], n_models=5)
ensemble.train_ensemble(X_train, y_train, epochs=100)

# ============================================
# EVALUATION AND COMPARISON
# ============================================

mc_dropout = MCDropoutPredictor(mc_model, n_samples=100)
mc_pred, mc_unc, mc_time = mc_dropout.predict_with_uncertainty(X_test)

ens_pred, ens_unc, ens_time = ensemble.predict_with_uncertainty(X_test)

mc_accuracy = np.mean(np.argmax(mc_pred, axis=1) == y_test.numpy())
ens_accuracy = np.mean(np.argmax(ens_pred, axis=1) == y_test.numpy())

print(f"{'Metric':<30} {'MC Dropout':<15} {'Deep Ensemble':<15}")
print("-" * 60)
print(f"{'Accuracy':<30} {mc_accuracy:.4f}{'':10} {ens_accuracy:.4f}")
print(f"{'Mean Uncertainty':<30} {np.mean(mc_unc):.4f}{'':10} {np.mean(ens_unc):.4f}")
print(f"{'Inference Time (s)':<30} {mc_time:.4f}{'':10} {ens_time:.4f}")
\end{lstlisting}

\begin{keyconcept}[Key Findings]
\begin{enumerate}
    \item Deep Ensemble typically achieves higher accuracy (model averaging)
    \item Both methods show higher uncertainty for incorrect predictions (good calibration)
    \item MC Dropout is slower at inference (100 passes vs.\ 5)
    \item For HITL: Set threshold, e.g., flag for review if uncertainty $> 0.25$
\end{enumerate}
\end{keyconcept}

\subsection{Exercise 3: A/B Testing Framework for Threshold Optimization}

\textbf{Task:} Design an A/B testing framework for optimizing human intervention thresholds.

\begin{lstlisting}[style=python, caption={Content moderation A/B testing framework}]
import numpy as np
import pandas as pd
from scipy import stats

class ContentModerationABTest:
    """
    A/B testing framework for content moderation thresholds.

    Demonstrates all five dimensions:
    1. Uncertainty Detection: AI toxicity scores
    2. Intervention Design: Human review queue
    3. Timing: Real-time moderation
    4. Stakes Calibration: Cost parameters
    5. Feedback Integration: Learning optimal threshold
    """

    def __init__(self):
        self.human_accuracy = 0.95

    def generate_content(self, n_samples=10000):
        """Generate synthetic content with toxicity scores."""
        true_labels = np.random.binomial(1, 0.15, n_samples)  # 15% toxic

        ai_scores = np.where(
            true_labels == 1,
            np.random.beta(8, 2, n_samples),
            np.random.beta(2, 8, n_samples)
        )
        ai_scores = np.clip(
            ai_scores + np.random.normal(0, 0.08, n_samples), 0, 1)

        return pd.DataFrame({
            'ai_score': ai_scores,
            'true_label': true_labels
        })

    def moderate(self, df, threshold):
        """Apply moderation with given threshold."""
        results = df.copy()

        # Human review for uncertain cases
        results['needs_review'] = (
            (results['ai_score'] > threshold) &
            (results['ai_score'] < (1 - threshold))
        )

        # AI-only decisions
        ai_only = ~results['needs_review']
        results.loc[ai_only, 'decision'] = (
            results.loc[ai_only, 'ai_score'] > 0.5).astype(int)

        # Human decisions (95% accurate)
        human = results['needs_review']
        n_human = human.sum()
        correct = np.random.binomial(1, self.human_accuracy, n_human)
        results.loc[human, 'decision'] = np.where(
            correct,
            results.loc[human, 'true_label'],
            1 - results.loc[human, 'true_label']
        )

        return results

    def calculate_metrics(self, results):
        """Calculate performance and cost metrics."""
        n = len(results)
        accuracy = (results['decision'] == results['true_label']).mean()

        tp = ((results['decision'] == 1) & (results['true_label'] == 1)).sum()
        fp = ((results['decision'] == 1) & (results['true_label'] == 0)).sum()
        fn = ((results['decision'] == 0) & (results['true_label'] == 1)).sum()

        precision = tp / (tp + fp) if (tp + fp) > 0 else 0
        recall = tp / (tp + fn) if (tp + fn) > 0 else 0
        review_rate = results['needs_review'].mean()

        # Cost analysis (Stakes Calibration)
        daily_volume = 100000
        daily_cost = (
            review_rate * daily_volume * 0.50 +     # $0.50 per review
            (fn / n) * daily_volume * 100.00 +      # $100 per missed toxic
            (fp / n) * daily_volume * 5.00          # $5 per false positive
        )

        return {
            'accuracy': accuracy,
            'precision': precision,
            'recall': recall,
            'review_rate': review_rate,
            'daily_cost': daily_cost
        }

    def run_test(self, threshold_a, threshold_b, n_samples=20000):
        """Run A/B test comparing two thresholds."""
        data_a = self.generate_content(n_samples)
        data_b = self.generate_content(n_samples)

        results_a = self.moderate(data_a, threshold_a)
        results_b = self.moderate(data_b, threshold_b)

        metrics_a = self.calculate_metrics(results_a)
        metrics_b = self.calculate_metrics(results_b)

        # Statistical test
        acc_a = (results_a['decision'] == results_a['true_label']).sum()
        acc_b = (results_b['decision'] == results_b['true_label']).sum()

        p_a, p_b = acc_a/n_samples, acc_b/n_samples
        p_pool = (acc_a + acc_b) / (2 * n_samples)
        se = np.sqrt(p_pool * (1-p_pool) * (2/n_samples))
        z = (p_b - p_a) / se if se > 0 else 0
        p_value = 2 * (1 - stats.norm.cdf(abs(z)))

        cost_diff = metrics_a['daily_cost'] - metrics_b['daily_cost']
        if p_value < 0.05 and (p_b - p_a) > 0.01:
            print(f"Recommendation: Use threshold B (higher accuracy)")
        elif cost_diff > 1000:
            print(f"Recommendation: Consider B (${cost_diff:,.0f}/day savings)")
        else:
            print("Recommendation: No clear winner, continue testing")

# Run demonstration
ab_test = ContentModerationABTest()
ab_test.run_test(0.25, 0.35)
ab_test.run_test(0.35, 0.45)
\end{lstlisting}

\subsection{Exercise 5: Nest Thermostat Redesign---Complete Solution}

\textbf{Task:} Redesign the Nest thermostat Home/Away detection using the five-dimension framework.

\begin{lstlisting}[style=python, caption={Complete ImprovedNestController implementation}]
from dataclasses import dataclass
from typing import Literal, Optional
from datetime import datetime, timedelta
import numpy as np

@dataclass
class PresenceSignals:
    motion_detected: bool
    time_since_motion: timedelta
    phone_in_geofence: bool
    phone_confidence: float
    smart_device_activity: bool
    typical_home_time: bool
    recent_manual_override: bool

@dataclass
class EnvironmentContext:
    outside_temp: float
    current_setpoint: float
    weather_forecast: str

@dataclass
class PresenceResult:
    status: Literal['HOME', 'AWAY', 'ASK_USER']
    confidence: float
    reasoning: str


class ImprovedNestController:
    """Redesigned Nest with five-dimension framework."""

    def __init__(self):
        self.false_away_count = 0

    # DIMENSION 1: Uncertainty Detection
    def calculate_presence_probability(
            self, signals: PresenceSignals) -> tuple:
        probability = 0.5
        confidence = 0.5

        if signals.motion_detected:
            probability += 0.35
            confidence += 0.2
        elif signals.time_since_motion < timedelta(minutes=30):
            probability += 0.15
        elif signals.time_since_motion > timedelta(hours=2):
            probability -= 0.35
            confidence += 0.15

        if signals.phone_in_geofence:
            probability += 0.30 * signals.phone_confidence
            confidence += 0.15 * signals.phone_confidence
        elif signals.phone_confidence > 0.8:
            probability -= 0.30
            confidence += 0.15

        if signals.smart_device_activity:
            probability += 0.15
        if signals.typical_home_time:
            probability += 0.10
        if signals.recent_manual_override:
            confidence -= 0.2
        if self.false_away_count > 3:
            probability += 0.1
            confidence -= 0.1

        return np.clip(probability, 0, 1), np.clip(confidence, 0.2, 1)

    # DIMENSION 4: Stakes Calibration
    def calculate_stakes(self, context: EnvironmentContext) -> str:
        if context.outside_temp < 35 or context.outside_temp > 95:
            return 'high'
        elif abs(context.outside_temp - context.current_setpoint) > 25:
            return 'high'
        elif abs(context.outside_temp - context.current_setpoint) > 15:
            return 'medium'
        return 'low'

    def get_threshold(self, stakes: str) -> float:
        return {'low': 0.70, 'medium': 0.85, 'high': 0.95}[stakes]

    # DIMENSION 2: Intervention Design
    def create_prompt(
            self, prob: float, context: EnvironmentContext) -> dict:
        weather_note = ""
        if context.outside_temp < 40:
            weather_note = f"It's {context.outside_temp}F outside."
        elif context.outside_temp > 85:
            weather_note = f"It's {context.outside_temp}F outside."

        return {
            'title': "Home Status Check",
            'message': f"Haven't detected activity. {weather_note}",
            'options': ["I'm home", "I'm out", "Ask later",
                        "Don't ask today"]
        }

    # MAIN DECISION LOGIC (Dimensions 1 + 3 + 4)
    def determine_presence(
            self, signals: PresenceSignals,
            context: EnvironmentContext) -> PresenceResult:
        probability, confidence = self.calculate_presence_probability(
            signals)
        stakes = self.calculate_stakes(context)
        threshold = self.get_threshold(stakes)

        if probability > threshold:
            return PresenceResult(
                'HOME', probability,
                f"High confidence home. Stakes: {stakes}")
        elif probability < (1 - threshold):
            return PresenceResult(
                'AWAY', 1-probability,
                f"High confidence away. Stakes: {stakes}")
        else:
            return PresenceResult(
                'ASK_USER', confidence,
                f"Uncertain. Stakes: {stakes}. Asking user.")

    # DIMENSION 5: Feedback Integration
    def process_feedback(self, prediction: str, response: str):
        if prediction == 'AWAY' and response == "I'm home":
            self.false_away_count += 1
            print("Learning: Was wrong. Now more conservative.")


# Demonstration
def demo():
    controller = ImprovedNestController()

    signals = PresenceSignals(
        motion_detected=False,
        time_since_motion=timedelta(hours=1),
        phone_in_geofence=False,
        phone_confidence=0.6,
        smart_device_activity=False,
        typical_home_time=True,
        recent_manual_override=False
    )
    context = EnvironmentContext(
        outside_temp=28,
        current_setpoint=70,
        weather_forecast='extreme_cold'
    )

    result = controller.determine_presence(signals, context)
    print(f"Status: {result.status}")
    print(f"Confidence: {result.confidence:.0%}")
    print(f"Reasoning: {result.reasoning}")

    if result.status == 'ASK_USER':
        prompt = controller.create_prompt(result.confidence, context)
        print(f"\nNotification: {prompt['message']}")
        print(f"Options: {prompt['options']}")

demo()
\end{lstlisting}

\paragraph{Expected Output.}
\begin{lstlisting}[language={}, caption={Expected output from demo()}]
Status: ASK_USER
Confidence: 45%
Reasoning: Uncertain. Stakes: high. Asking user.

Notification: Haven't detected activity. It's 28F outside.
Options: ["I'm home", "I'm out", "Ask later", "Don't ask today"]
\end{lstlisting}

% ============================================================================
\section{Key Takeaways}
% ============================================================================

\subsection{Five-Dimension Framework Summary}

\begin{table}[htbp]
\caption{Five-dimension framework: implementation patterns}
\label{tab:framework-patterns}
\centering
\begin{tabular}{@{}p{3.5cm}p{8cm}@{}}
\toprule
\textbf{Dimension} & \textbf{Implementation Pattern} \\
\midrule
Uncertainty Detection & \code{confidence = model.predict\_uncertainty(x)} \\
Intervention Design & Clear prompt + multiple response options \\
Timing & \code{delay = base\_delay * stakes\_multiplier} \\
Stakes Calibration & \code{threshold = cost\_FP / (cost\_FP + cost\_FN)} \\
Feedback Integration & \code{model.update(user\_correction)} \\
\bottomrule
\end{tabular}
\end{table}

\subsection{General HITL Decision Pattern}

\begin{lstlisting}[style=python, caption={General HITL decision pattern}]
def should_intervene(confidence, stakes):
    threshold = get_threshold(stakes)

    if confidence > threshold:
        return "act_automatically"
    elif confidence < (1 - threshold):
        return "act_automatically_opposite"
    else:
        return "ask_human"
\end{lstlisting}

\bigskip
\begin{center}
\rule{0.5\textwidth}{0.4pt}
\end{center}
\bigskip

\noindent\textit{These solutions demonstrate practical application of HITL principles from Chapter~1 and Technical Appendix~1A. All answers are framework-aligned to provide instructors with confidence in facilitating discussions, consistency in grading across sections, and depth to handle advanced student questions.}
